\subsection{Предсказывает ли качество карты безопасность (H4)}

Для каждой сцены сопоставлены mIoU карты при истинных позах и доля нарушений планировщика без калибровки (рис.~\ref{fig:h4}). Ранговая корреляция близка к нулю: $\rho = -0{,}16$ ($p = 0{,}57$, 15 сцен); для IoU только запретных классов знак даже обратный ожидаемому ($\rho = +0{,}51$, $p = 0{,}05$). При разбросе mIoU сцен от 0{,}17 до 0{,}34 метрика карты не говорит, насколько опасно по ней ехать: безопасность определяется наихудшим недопокрытием объекта, а mIoU усредняет по клеткам. Гипотеза H4 подтверждается (с оговоркой о малом числе сцен).

\begin{figure}[H]
\centering
\includegraphics[width=0.78\linewidth]{figures/miou_vs_violations_pb.png}
\caption{Качество карты против доли нарушений планировщика без калибровки (одна точка --- сцена, L0)}
\label{fig:h4}
\end{figure}

\subsection{Планировщик на ужесточённой карте (H5)}\label{sec:bridge}

На уровне L2 запретные зоны совместной ветки --- те же, что в миссиях, --- масштабированы коэффициентом $m$ (угловая зона~\cite{dovgopolik2025}); на задачах, решаемых точным поиском, сравниваются три выборочных планировщика с бюджетом 1500 расширений (табл.~\ref{tab:bridge}, рис.~\ref{fig:bridge}): Р --- RRT-Connect с равномерной выборкой, У --- угловая динамическая зона по опубликованному правилу (расширение конуса, пока образцы попадают в препятствие), А --- предложенная модификация: конус расширяется после каждого заблокированного расширения дерева и сужается после успешного.

\begin{table}[H]
\centering\scriptsize
\caption{Выборочные планировщики на ужесточённых картах (растение и велосипед, L2, 233 задачи, 2 запуска на задачу)}
\label{tab:bridge}
\begin{tabular}{ccc|ccc|ccc|ccc}
\toprule
$m$ & своб. & решаемо & \multicolumn{3}{c|}{решено за бюджет} & \multicolumn{3}{c|}{итераций (медиана)} & \multicolumn{3}{c}{длина / кратчайший} \\
 & & & Р & У & А & Р & У & А & Р & У & А \\
\midrule
0{,}0 & 0{,}96 & 0{,}97 & 0{,}93 & 0{,}91 & 0{,}93 & 66 & 192 & 94 & 1{,}15 & 1{,}14 & 1{,}11 \\
0{,}5 & 0{,}89 & 0{,}67 & 0{,}75 & 0{,}71 & 0{,}73 & 88 & 236 & 113 & 1{,}18 & 1{,}16 & 1{,}13 \\
1{,}0 & 0{,}79 & 0{,}30 & 0{,}89 & 0{,}89 & 0{,}86 & 92 & 213 & 111 & 1{,}19 & 1{,}17 & 1{,}17 \\
1{,}5 & 0{,}70 & 0{,}16 & 0{,}85 & 0{,}80 & 0{,}81 & 84 & 256 & 99 & 1{,}21 & 1{,}18 & 1{,}14 \\
\bottomrule
\end{tabular}

\end{table}

\begin{figure}[H]
\centering
\includegraphics[width=\linewidth]{figures/planner_bridge_pb.png}
\caption{Слева: доля свободного пространства и решаемых задач; в центре: доля решённых за бюджет $B$ при откалиброванном запасе; справа: длина первого найденного пути относительно кратчайшего}
\label{fig:bridge}
\end{figure}

\begin{itemize}[nosep]
  \item Ограничивает свободное пространство, а не поиск: при откалиброванном запасе решаемы лишь $30\,\%$ задач (при $m = 0$ --- $97\,\%$), а решаемые задачи все три планировщика решают за бюджет примерно одинаково часто ($0{,}86$--$0{,}89$).
  \item Опубликованное правило угловой зоны требует в 2{,}3--3{,}0 раза больше расширений, чем равномерная выборка: на ужесточённых картах проходы часто смещены от прямой между деревьями, а конус расширяется только при образцах в препятствии, которых в почти свободном пространстве мало.
  \item Модификация с расширением конуса по заблокированным расширениям сокращает это до 1{,}2--1{,}4 раза. Её первые пути не длиннее, чем у других: $1{,}11$--$1{,}17$ от кратчайшего против $1{,}15$--$1{,}21$ у равномерной выборки; при откалиброванном запасе она наравне с опубликованным правилом ($1{,}17$). Т.~е. свойство квазиоптимальности исходного метода сохраняется.
\end{itemize}
Это первый, предварительный результат для ВКР: вклад планировщика на картах с калиброванным запасом --- качество и стоимость пути в суженном пространстве; адаптация угловой зоны к смещённым проходам --- конкретная исследовательская задача.
